% TOP_PROBLEM: {"id":30007546,"problem_number":"LOCAL-30007546","title":"Declining business dynamism","category_id":21,"category":{"id":21,"name":"economics","display_name":"Economics","description":"Open economic theory statements, published research agendas, and proposed scoped specifications; question provenance and historical priority are qualified per record.","slug":"economics","order_index":21,"created_at":"2026-10-08T00:00:00Z"},"status":"research_question","external_url":null,"legacy_ids":[],"published":false,"statement":"Identify a structural decomposition of declining business dynamism that jointly fits aggregate trends and industry-level evidence.","economics":{"original_id":35,"field":"Growth and productivity","kind":"identification","formalization_basis":"adapted_research_question","source_status":"explicit_open","priority_status":"earliest_found","priority_note":"This is the earliest explicit relevant statement verified in this audit, not a claim of original priority; older literature and earlier versions may contain antecedents.","earliest_verified_reference":{"authors":["Brian C. Albrecht","Ryan A. Decker"],"title":"Rising Markups and Declining Business Dynamism: Evidence From the Industry Cross Section","year":2024,"url":"https://www.federalreserve.gov/econres/notes/feds-notes/rising-markups-and-declining-business-dynamism-evidence-from-the-industry-cross-section-20240308.html","doi":null,"locator":"Section 2, “Related literature,” paragraphs on the absence of an overriding explanation and reconciling time-series with cross-sectional evidence","evidence_summary":"The authors say no single explanation dominates the dynamism literature and identify the lack of a framework reconciling aggregate trends with industry evidence. Their analysis cautions against a simple markup-based explanation."},"current_reference":{"authors":["Brian C. Albrecht","Ryan A. Decker"],"title":"Rising Markups and Declining Business Dynamism: Evidence From the Industry Cross Section","year":2024,"url":"https://www.federalreserve.gov/econres/notes/feds-notes/rising-markups-and-declining-business-dynamism-evidence-from-the-industry-cross-section-20240308.html","doi":null,"locator":"Section 2, “Related literature,” paragraphs on the absence of an overriding explanation and reconciling time-series with cross-sectional evidence","evidence_summary":"The authors say no single explanation dominates the dynamism literature and identify the lack of a framework reconciling aggregate trends with industry evidence. Their analysis cautions against a simple markup-based explanation."},"formalization_note":"The source explicitly identifies the lack of one overriding explanation and of a reconciliatory framework. The parameter-block decomposition is adapted."},"statusQualification":"Published question or research agenda verified at the cited locator. The mathematical specification is adapted for this catalogue and includes additional assumptions; source publication does not establish first-ever priority or certify this exact model remains unsolved. The source explicitly identifies the lack of one overriding explanation and of a reconciliatory framework. The parameter-block decomposition is adapted.","statusReviewedAt":"2026-10-08"}
% FIRST_POSTED: 2026-10-08
% ENTRY_KIND: statement_only
% !TeX program = lualatex
\documentclass[11pt]{article}
\usepackage{fontspec}
\usepackage[a4paper,margin=25mm]{geometry}
\usepackage{amsmath,amssymb,mathtools}
\usepackage{xurl}
\usepackage[hidelinks]{hyperref}
\newcommand{\sref}[2]{\hyperlink{source:#1:#2}{[S#2]}}
\setlength{\emergencystretch}{3em}
\begin{document}
% problemId:30007546
\section*{Declining business dynamism}\label{30007546}
\subsection{Definitions and mathematical statement}
Let \(D_t\) be a predeclared business-dynamism index built from firm entry, exit, and employment reallocation in a fixed population. Consider a structural firm-dynamics model with parameter blocks \(\theta=(\theta^N,\theta^C,\theta^I,\theta^R)\), representing demographics and labor supply, competitive conditions, intangible technology and diffusion, and regulation. Specify how each block affects entry costs, incumbent innovation, and firm growth; a residual trend must not silently receive a causal label.
For initial and final parameter vectors \(\theta_0,\theta_1\), let \(D(\theta)\) be the equilibrium index. Decompose \(D(\theta_1)-D(\theta_0)\) into block contributions using a declared Shapley ordering average over counterfactual combinations. Identify those contributions from moments containing both aggregate trends and cross-industry variation, augmented by credible policy or technological instruments. In particular, test whether the same parameter changes explain industry-level relationships among markups, entry, and reallocation rather than only matching aggregate comovement.
Report observationally equivalent decompositions and parameter uncertainty, distinguish measured markups from primitive market power, and assess robustness to firm-population coverage. A successful result reconciles time-series decline with cross-sectional evidence and predicts held-out changes. The target is a quantitatively disciplined explanation for a stated period and economy, not a presumption that a single driver must dominate universally.

\par\textbf{Formulation and scope.} The source explicitly identifies the lack of one overriding explanation and of a reconciliatory framework. The parameter-block decomposition is adapted.

\par\textbf{Open-question provenance.} An explicit question or research agenda was verified at the source locator. This does not by itself certify that every later version remains unresolved \sref{30007546}{1}.

\par\textbf{Historical priority.} This is the earliest explicit relevant statement verified in this audit, not a claim of original priority; older literature and earlier versions may contain antecedents.

\subsection{Short English statement}
Identify a structural decomposition of declining business dynamism that jointly fits aggregate trends and industry-level evidence.

\subsection{Sources}
\begin{itemize}\raggedright
\item[\textnormal{[S1]}]\hypertarget{source:30007546:1}{} Brian C. Albrecht, Ryan A. Decker. 2024. Rising Markups and Declining Business Dynamism: Evidence From the Industry Cross Section. Section 2, “Related literature,” paragraphs on the absence of an overriding explanation and reconciling time-series with cross-sectional evidence.
\url{https://www.federalreserve.gov/econres/notes/feds-notes/rising-markups-and-declining-business-dynamism-evidence-from-the-industry-cross-section-20240308.html}.
{\small\textit{Earliest verified question/agenda reference; historical priority qualified below. The authors say no single explanation dominates the dynamism literature and identify the lack of a framework reconciling aggregate trends with industry evidence. Their analysis cautions against a simple markup-based explanation.}}
\end{itemize}
\end{document}
